\chapter{Conclusiones}\label{cap:conclusiones}

El capítulo se organiza del siguiente modo.
En primer lugar se responde a las cuatro preguntas de investigación, y a continuación se evalúa cada contribución frente a los objetivos de los que sale.
Después se agrupan las amenazas a la validez y se describe el paquete de replicación.
Por último se recoge el trabajo futuro y se cierra con las conclusiones técnicas del trabajo.

\section{Respuesta a las preguntas de investigación}\label{sec:conclusiones:rq}

Cada pregunta se reproduce con el enunciado de la Sección~\ref{sec:introduccion:objetivos}, y la respuesta remite al apartado que la sostiene.

\noindent\textbf{RQ1.} \emph{¿Qué APIs de correo deben incluirse en un análisis representativo del dominio, y bajo qué criterios reproducibles?}

El catálogo reúne ocho APIs, y entre ellas aparece al menos una vez cada una de las cuatro formas de cobro que fija el primer criterio de la consulta (Sección~\ref{sec:propuesta:seleccion}).
Dos de esas formas son planes, por volumen mensual o con un límite diario adicional, y las otras dos son el pago por uso puro y el producto vendido como añadido de otro.
La Sección~\ref{sec:propuesta:respuesta} cierra la RQ1 con la condición completa: el conjunto debe cubrir esas cuatro formas, y sus APIs deben diferir en el rasgo diferenciador de límites que pide el cuarto criterio (Sección~\ref{sec:propuesta:criterios}).
Ese rasgo es, por ejemplo, un \textit{rate} publicado por \textit{endpoint}, un caudal negociado no publicado o un límite diario.
Cada API se documentó con la trazabilidad que exige la Sección~\ref{sec:propuesta:trazabilidad}.
Ninguna API se sustituyó, así que la lista de reserva (Apartado~\ref{sec:propuesta:reserva}) no llegó a usarse.
El Capítulo~\ref{cap:hallazgos} analiza las mismas ocho de las que parte la selección.

Los criterios son reproducibles, pero no en la misma medida.
La cobertura se comprueba contra las páginas de precios de cada proveedor.
La consulta al modelo de lenguaje que completó la lista está en el Anexo~\ref{anx:uso-ia}, que reproduce el prompt literal y la respuesta íntegra, con una ficha que declara además lo que no consta.
El repositorio de datos que acompaña al trabajo conserva el registro original.
Constan también la fecha de la consulta y el modelo usado, pero repetir la consulta no garantiza la misma respuesta.

La respuesta tiene límites de dos clases.
La primera clase está en las fuentes: la primera del criterio, las APIs de un Trabajo Fin de Máster previo del mismo programa, no pudo ejecutarse, y la consulta la sustituyó.
La segunda clase son tres premisas de la respuesta que la construcción corrigió en la ventana de toma de datos (Apartado~\ref{sec:propuesta:ventana}).
Amazon SES entró en el catálogo el 1 de septiembre de 2026 como pago por uso puro, y el proveedor vendía planes desde el 21 de julio de 2026 (Apartado~\ref{subsec:hallazgos:ses}).
Su plan à la carte no cobra suscripción y cobra por correo (Tabla~\ref{tab:hallazgos:ses}), y de ahí se infiere que la forma sigue representada.
La respuesta presentaba además la combinación de cuota diaria y mensual como aportación de Brevo, y esa combinación ya estaba en el plan Free de Resend.
La Sección~\ref{sec:propuesta:respuesta} recoge además otros dos límites de la primera clase: los cinco criterios solo decidieron cuatro de las ocho APIs, y el tamaño del catálogo no se deriva de la cobertura.

La tercera premisa es la de Mailgun.
La respuesta registrada la eligió como el caso limpio de \textit{rate} documentado, pero su documentación de desarrollador no cifra el umbral (Sección~\ref{sec:evolucion:inventario}).
El único caudal que el proveedor publica para su API REST, 1.500 peticiones por segundo, está en sus preguntas frecuentes y acotado a una IP de cliente.
Esa premisa es la que justificó su entrada, porque la respuesta le atribuye la misma forma de cobro que a Postmark, SendGrid y MailerSend.
La construcción encontró otro rasgo de límites: Mailgun es la única del catálogo que contesta por escrito sobre la concurrencia de su API REST (Sección~\ref{sec:propuesta:respuesta}).

\noindent\textbf{RQ2.} \emph{¿Es ACM2YAML suficientemente expresivo para representar la capacidad contractual de las APIs de correo? ¿Qué extensiones requiere?}

Con las ocho APIs del catálogo, la respuesta a la primera pregunta es que no.
La respuesta a la segunda son las nueve propuestas de la Tabla~\ref{tab:propuestas}, siete de expresividad y dos de contrato, más una rectificación de la guía (Sección~\ref{sec:evolucion:respuesta}).
Las carencias de P-01, P-02 y P-03 están en dimensiones que v0.3 ya modela: el esquema las declara con valores escalares mientras el dominio las vende por listas, las consume por productos o las acota en otra métrica.
P-06 añade el estado de la suscripción y la reputación, un eje que el esquema solo sabe declarar dentro de cada \textit{endpoint} y repetido en todos, y que aparece en seis de las ocho APIs.
P-07 y P-09 tratan el exceso sobre la cuota, y P-08 pide una cuota que cuente entidades distintas y que no se reponga.
Los dos defectos de contrato, P-04 y P-05, son de otro tipo: el esquema admite el valor y el consumidor de referencia falla al leerlo.
El dominio pide más de lo que llega a propuesta.
La Tabla~\ref{tab:sinforma} recoge otros 12 elementos promovidos con la forma esbozada y sin fijar, y la Sección~\ref{sec:evolucion:fuera} deja fuera otros cuatro con su motivo.
Si esta respuesta vale fuera del catálogo se discute en la Sección~\ref{sec:conclusiones:externa}.

\noindent\textbf{RQ3.} \emph{¿Cómo se relacionan los límites de uso y los \textit{rates} en las APIs de correo, y qué implicaciones tiene esa relación para el consumidor?}

La cuota es la restricción activa cuando el \textit{rate} y la cuota tienen valor en la Datasheet, publicado o medido, y la cuenta está sana (Sección~\ref{sec:hallazgos:limites-rates}).
Eso ocurre en seis de las ocho APIs.
En su ventana, el \textit{rate} admite entre 2,6 veces la cuota (MailerSend Professional 2M) y 129.600.000 veces (el Free de Postmark, con el \textit{rate} en ámbar).
Las otras dos APIs quedan fuera de la regla.
SES no publica la cuota ni el caudal de producción, así que la relación no se puede calcular antes de sacar la cuenta del entorno de pruebas.
En Mailchimp Transactional la relación se invierte: en la cuenta de prueba, la franquicia de 500 correos con un tope de 25 por hora tarda 20 horas en agotarse, y el caudal es lo único que fija el tiempo.

Para el consumidor, la relación tiene tres implicaciones.
En esas seis APIs puede dimensionar el volumen sin mirar el \textit{rate}, siempre que la cuenta esté sana.
El estado de la suscripción y la reputación explican más de la capacidad observada que el plan, y en tres APIs el sondeo midió la capacidad de un estado y no la del plan.
Y solo en la mitad del catálogo, la que ofrece un endpoint de introspección, puede verificar la relación antes de enviar.

\noindent\textbf{RQ4.} \emph{¿Cómo pueden compararse las APIs de correo de forma sistemática para responder a preguntas de decisión que involucran simultáneamente requisitos funcionales y no funcionales?}

Las APIs se comparan con un catálogo de 32 preguntas de decisión (29 no funcionales y 3 funcionales) construido sobre la gramática de Pricing Intelligence \citep{cavero2026pricingintelligence}.
El procedimiento resuelve cada pregunta API a API y registra la fuente de cada respuesta (Secciones~\ref{sec:hallazgos:catalogo-pi} y~\ref{sec:hallazgos:comparacion}).
Por su clase, el instrumental del canal API alcanza 25 de las 29 no funcionales.
Las preguntas de HARVEY4API, 4 de esas 25, se clasifican por especificación, porque la herramienta no se ha ejecutado.
Sobre cada Datasheet, la cobertura va de 16 preguntas en Brevo a 22 en Resend (Sección~\ref{sec:hallazgos:comparacion}).
La matriz de \textit{features} por plan responde dos de las tres funcionales, fuera del esquema, y la tercera, la retención, espera sitio en él.
La parte funcional se resuelve por lectura porque el canal API que fija el alcance del Capítulo~\ref{cap:catalogo} no tiene instrumento para ella.
Una de las funcionales mezcla un requisito funcional con uno no funcional: pregunta qué API da IP dedicada al menor precio para 200.000 correos al mes.
Se responde componiendo a mano la matriz y la Datasheet (Apartado~\ref{subsec:hallazgos:funcional}), y responderla con instrumento exigiría el vínculo entre el iPricing del canal web y la Datasheet del canal API.
En 12 preguntas hay respuesta para las ocho APIs a la vez.

\section{Evaluación de las contribuciones frente a los objetivos}\label{sec:conclusiones:evaluacion}

La Tabla~\ref{tab:conclusiones:evaluacion} da el estado de cada objetivo y de la contribución que le corresponde, con las secciones de la memoria que lo sostienen.
Los párrafos que siguen justifican cada estado.
Empiezan por OBJ-03 y OBJ-04, donde el Capítulo~\ref{cap:introduccion} sitúa el componente de investigación.

\begin{table}[H]
\centering
\small
\setlength{\tabcolsep}{3pt}
\begin{tabular}{|l|c|c|p{1.9cm}|p{2.3cm}|p{4.6cm}|}
\hline
\textbf{Objetivo} & \textbf{RQ} & \textbf{Contrib.} & \textbf{Estado del objetivo} & \textbf{Estado de la contribución} & \textbf{Secciones que lo sostienen} \\
\hline
OBJ-01 & RQ1 & C-01 & cumplido en parte & cumplida en parte & \ref{sec:propuesta:seleccion}, \ref{sec:hallazgos:apis} \\
\hline
OBJ-02 & RQ2--RQ4 & C-01 & cumplido & cumplida en parte & \ref{sec:propuesta:trazabilidad}, \ref{sec:evolucion:metodo}, \ref{sec:evolucion:contrato} \\
\hline
OBJ-03 & RQ2 & C-02 & cumplido en parte & cumplida & \ref{sec:evolucion:inventario}, \ref{sec:evolucion:fuera}, \ref{sec:evolucion:respuesta} \\
\hline
OBJ-04 & RQ3 & C-03 & cumplido & cumplida & \ref{sec:hallazgos:limites-rates}, \ref{sec:hallazgos:conclusiones} \\
\hline
OBJ-05 & RQ4 & C-04 & cumplido & cumplida & \ref{sec:hallazgos:catalogo-pi}, \ref{sec:hallazgos:comparacion} \\
\hline
\end{tabular}
\caption{Estado de cada objetivo específico y de la contribución que le corresponde. RQ: la pregunta de investigación de la que se deriva el objetivo. OBJ-02 sirve de medio a RQ2, RQ3 y RQ4, porque el Capítulo~\ref{cap:introduccion} presenta el catálogo como condición necesaria. C-01 aparece dos veces porque responde a OBJ-01 y a OBJ-02. Los párrafos de esta sección justifican cada estado.}
\label{tab:conclusiones:evaluacion}
\end{table}

\textbf{OBJ-03 y C-02.} El criterio de promoción de la Sección~\ref{sec:propuesta:conformidad}, aplicado en el Capítulo~\ref{cap:evolucion}, promueve 24 de los 35 elementos no soportados del inventario (Sección~\ref{sec:evolucion:inventario}).
Las nueve propuestas de la Tabla~\ref{tab:propuestas} dan forma a ocho de esos promovidos, y la novena, P-05, no sale de ningún elemento del inventario.
La Tabla~\ref{tab:sinforma} recoge otros 12 promovidos con una forma esbozada y sin fijar.
Los cuatro promovidos restantes quedan fuera por los motivos de la Sección~\ref{sec:evolucion:fuera}.
OBJ-03 queda cumplido en parte porque pide la extensión que corresponde a cada carencia identificada, y los promovidos de la Tabla~\ref{tab:sinforma} no tienen todavía la forma fijada.
De las nueve, siete son de expresividad y dos de contrato, P-04 y P-05.
C-02 se cumple: cada una de las siete de expresividad responde a un caso real que v0.3 no representa, con su prueba de recurrencia o de traza.
Las dos de contrato quedan fuera del enunciado de C-02, porque v0.3 sí representa el valor y es el consumidor de referencia el que falla.
La rectificación corrige la guía y no extiende el esquema.

\textbf{OBJ-04 y C-03.} La Sección~\ref{sec:hallazgos:limites-rates} caracteriza la relación entre límites y \textit{rates} en las ocho APIs.
En seis de ellas la cuota es la restricción activa, y las Tablas~\ref{tab:limites-rates} y~\ref{tab:limites-rates-b} lo muestran con el cociente entre \textit{rate} y cuota de cada plan.
La lectura geométrica de la capacidad \citep{ruizcortes2025cajanegra} se aplica con las curvas $C(t)$ que PRIME4API dibuja en local para cuatro APIs (Figura~\ref{fig:capacidad}).
En Mailgun, con dos contadores sobre la misma ventana, y en Mailchimp Transactional, con una cuota sin ventana, la capacidad no cabe en una curva de ese tipo (Sección~\ref{sec:hallazgos:limites-rates}).
La caracterización responde a la tercera observación del Capítulo~\ref{cap:estadoDelArte}: en seis APIs la cuota enmascara al \textit{rate}, y en Mailchimp Transactional el \textit{rate} enmascara a la cuota.
OBJ-04 y C-03 se cumplen, con las salvedades de la Sección~\ref{sec:conclusiones:amenazas}.

\textbf{C-01 frente a OBJ-01 y OBJ-02.} OBJ-02 se cumple, con la salvedad de la Sección~\ref{sec:conclusiones:constructo}.
Todas las Datasheets del catálogo siguen ACM2YAML v0.3, y cada valor lleva la URL de procedencia y la fecha de acceso, junto a la cita literal que lo sostiene (Sección~\ref{sec:propuesta:trazabilidad}).
Superan además el validador de formato de este trabajo, aunque usan dos formas que la guía no define: la unidad \texttt{MBs} y la lista \texttt{not\_found} (Sección~\ref{sec:evolucion:metodo}).
OBJ-01 se cumple en parte.
La primera fuente del criterio no se ejecutó, y la premisa por la que entró SES no se sostuvo.
Un tercero tampoco puede repetir la lista entera, que es lo que OBJ-01 le pide, porque la consulta registrada no garantiza la misma respuesta.
Postmark, además, no entra por ninguna de las tres fuentes: es el piloto con el que se calibró el método (Sección~\ref{sec:propuesta:piloto}).
C-01 se cumple en parte.
El catálogo sale revisado del ciclo de cada API, está versionado y cada dato remite a su fuente.
La contribución a SPHERE está preparada \citep{sphere}: los ocho pricings en Pricing2YAML pasan el analizador que usa la plataforma, y las 20 limitaciones que Pricing2YAML deja fuera quedan documentadas.
La carga por la interfaz de SPHERE no se ha probado, y la publicación depende de los responsables de la plataforma.
Que sea reutilizable por el resto de herramientas del ecosistema, la otra promesa de C-01, se sostiene en parte.
PRIME4API lee en local las ocho Datasheets, pero ninguna es plenamente consumible (Tabla~\ref{tab:verificacion}).
En MailerSend, Amazon SES y Mailgun falla una sola comprobación, y el Capítulo~\ref{cap:evolucion} no deriva de ese fallo ninguna propuesta de esquema.
HARVEY4API no se ha ejecutado.

\textbf{OBJ-05 y C-04.} El catálogo de preguntas existe, y el procedimiento se ha repetido sobre las ocho APIs (Sección~\ref{sec:hallazgos:comparacion}).
OBJ-05 y C-04 se cumplen con tres límites declarados.
Las preguntas de HARVEY4API se clasifican por especificación, y las preguntas de PRIME4API de Mailchimp Transactional están calculadas a mano.
Dos de las funcionales se responden leyendo la matriz de \textit{features}, y la tercera, la retención, espera sitio en el esquema.
El vínculo con el iPricing que daría instrumento a la pregunta funcional mixta queda como trabajo futuro.

\section{Amenazas a la validez}\label{sec:conclusiones:amenazas}

Cada apartado del Capítulo~\ref{cap:hallazgos} declara sus límites al cerrar, y el Capítulo~\ref{cap:evolucion} hace lo mismo con los suyos.
Esta sección los agrupa en cuatro ejes de validez, uno por apartado.
Cada amenaza parte de un caso real y dice qué sesga.
Después dice qué la mitiga en este trabajo y qué queda abierto.

\subsection{Validez de constructo}\label{sec:conclusiones:constructo}

La validez de constructo pregunta si lo medido es lo que se dice medir.
El sondeo midió en tres APIs un estado de la suscripción y no el plan.
En MailerSend era una cuenta Sandbox sin aprobar.
En SES era el entorno de pruebas, y en Mailchimp Transactional una cuenta de prueba.
En Brevo y Mailgun la cuenta era del plan gratuito.
Brevo describe, sin cifra, un estado de la suscripción pendiente de aprobación, y Mailgun un arranque de cuenta nueva cuyo caudal sube con la reputación.
Lo medido vale para el plan gratuito en el estado de esa cuenta, y la Datasheet de Mailgun referencia el caudal medido desde sus cuatro planes.
El sesgo está en la capacidad observada, que es la del estado.
MailerSend publica 120 peticiones por minuto para una cuenta sana, y la cuenta Sandbox sin aprobar del sondeo devolvió 10, la cifra que el proveedor publica para la reputación bajo revisión.
Lo mitiga que las tres Datasheets declaren el estado medido, y que la de MailerSend conserve aparte el valor documental del plan.
Queda abierto medir en cuentas de producción, y salir de esos estados exige alta de cuenta o gasto, que este trabajo no decide.

En Postmark y Mailgun el \textit{rate} se adoptó desde una cabecera, sin llegar al rechazo.
El sondeo de Postmark envió 200 peticiones sin obtener un 429, y en Mailgun alcanzarlo exige 80.000 peticiones en 10 segundos, una carga para la que el sondeo no se diseñó.
En SendGrid, el caudal de envío se publica en una página de preguntas frecuentes y ninguna cabecera lo deja observar, así que solo se contrastaría con una campaña de carga.
El sesgo está en la cifra: la de Postmark y Mailgun es la que anuncia la cabecera, la de SendGrid es la de la página, y en ninguna consta la que el proveedor aplica al rechazar.
Lo mitiga que las cifras de Postmark y Mailgun vayan en ámbar en su Datasheet.
La respuesta a RQ3 no cambia mientras el \textit{rate} aplicado siga por encima de lo que la cuota deja enviar en su ventana.
Es inferencia y no medida.
Queda abierto el cuarto peldaño de la escalera de validación, la medición formal \citep{datasheetusagetesting}.

La ventana de las cuotas tiene cuatro APIs en rojo, cuatro en ámbar y ninguna en verde para todas sus cuotas, y la ventana del caudal no la declara ninguna API.
La unidad de la cuota tiene respaldo en cuatro APIs, y en dos de ellas, Brevo y SES, lo aporta un sondeo.
En las otras cuatro no consta que un destinatario cuente como un correo a efectos de contador.
El sesgo alcanza a los dos pasos del cálculo de RQ3: el cociente entre \textit{rate} y cuota se toma sobre la ventana de la cuota, y los correos salen de multiplicar peticiones por destinatarios.
Lo mitiga que el semáforo de confianza dé color a la ventana de cada cuota.
El respaldo de la unidad lo declara API a API el Capítulo~\ref{cap:hallazgos}, así que ninguno de los dos datos se presenta como publicado sin serlo.
Queda abierto medir la unidad en Postmark, MailerSend, Mailgun y Mailchimp Transactional, las cuatro sin respaldo.

La trazabilidad de cada valor descansa en su cita literal, y la cita solo se comprueba de forma mecánica si se archivó el volcado de su página.
En MailerSend, la última comprobación registrada del verificador de citas deja 70 citas sin volcado en la carpeta contra la que se contrastan.
La mayoría son de páginas leídas sin guardar su volcado, y el resto son literales del sondeo, cuyo volcado es de otra fecha.
Postmark no tiene ningún volcado archivado, así que ninguna de sus citas se puede comprobar de forma mecánica (Sección~\ref{sec:conclusiones:replicacion}).
El sesgo está en la comprobación: las de páginas sin volcado no se pueden contrastar con la página tal como era el día de acceso.
Lo mitiga que la URL y la fecha de acceso consten en cada valor.
Queda abierto releer esas páginas y archivar su volcado.

Lo que este trabajo llama carencia de v0.3 depende de con qué se contrasta la Datasheet, y no existe validador oficial de Datasheets (Sección~\ref{sec:evolucion:metodo}).
La conformidad se comprueba con el validador de formato de este trabajo y con la verificación por uso sobre PRIME4API.
Las dos formas que la guía no define, \texttt{MBs} y \texttt{not\_found}, están declaradas.
Un validador propio corre el riesgo de medir contra su propia lectura de la guía.
Lo mitiga, en los defectos de contrato y en los elementos no soportados que entran por traza, usar como consumidor el motor del ecosistema, que no es de este trabajo.
Las carencias que entran solo por recurrencia, las de P-01, P-02, P-03 y P-06, descansan en la lectura de la guía que hace este trabajo.
P-06 muestra el riesgo.
La primera pasada, hecha sobre la guía, dio por fuera del mecanismo el cambio de carga por petición, y el código lo admitía.
La ausencia de mecanismo se establece sobre el texto de la guía, y en P-06 se contrastó también el código de referencia (Sección~\ref{sec:evolucion:compatibilidad}).
El rastreo de las carencias ya reportadas es del 10 de agosto de 2026.
Queda abierto contrastar las carencias con el código de la guía y con lo reportado después del 10 de agosto.

El catálogo de preguntas de decisión es de este trabajo, y la cobertura del instrumento depende de qué preguntas incluye.
El sesgo tiene dirección: quien escribe las preguntas conoce el instrumental, y eso puede subir la cobertura.
Es inferencia, pero basta para no leer la cobertura de este catálogo, un 78~\%, contra el 54~\% que el paper de Pricing Intelligence \citep{cavero2026pricingintelligence} obtiene sobre su banco de 56 plantillas.
Lo mitiga que las preguntas sigan su gramática, y que ni sus clases ni su denominador se muevan al pasar de cuatro APIs a ocho.
Queda abierto declarar el criterio con el que se generó el catálogo y por qué se detiene en 32 preguntas.

\subsection{Validez interna}\label{sec:conclusiones:interna}

La validez interna pregunta si los resultados dependen de las Datasheets o del montaje con que se obtienen.
La herramienta de verificación clona PRIME4API sin fijar la revisión, así que el resultado de la Tabla~\ref{tab:verificacion} depende de la revisión que se clone (Sección~\ref{sec:evolucion:metodo}).
La copia local del motor se clonó el 4 de septiembre de 2026, y su último cambio es del 18 de agosto.
Un cambio del motor entre dos ejecuciones puede cambiar el número de errores sin que cambie la Datasheet.
Las ocho ejecuciones no son simultáneas: Mailchimp Transactional y Mailgun se sometieron a la herramienta el 24 de septiembre de 2026, después de las otras seis.
Las de Mailchimp Transactional y Mailgun registran la misma revisión \texttt{8e6011c} que las de Postmark, Resend, SendGrid y MailerSend.
Lo mitiga, en parte, que esas seis ejecuciones con revisión registrada puedan repetirse contra ella.
Para las de Brevo y SES no hay mitigación, porque su revisión no consta.
Queda abierto fijar la revisión en la herramienta y registrarla en cada ejecución.

El Capítulo~\ref{cap:evolucion} recoge tres observaciones que son límites del propio método de sondeo (Sección~\ref{sec:evolucion:observaciones}).
En MailerSend el contador expuesto se actualiza por lotes y con retardo, y qué lo pone al día no se midió (O-11).
Un sondeo no puede usar el contador de la respuesta siguiente como coste de la anterior.
En Brevo, un destinatario que no se puede entregar no asegura el rechazo, porque con la cabecera de modo \textit{sandbox} el envío devolvió \texttt{201} (O-20).
En Mailchimp Transactional, desde 11 destinatarios o con un adjunto (el menor probado, de 5 MB), el envío queda en cola, y deja de existir la vía de medir por rechazo (O-22).
Las tres observaciones pueden hacer pasar por coste o por límite una señal que no lo es.
En Brevo lo mitiga un sondeo previo que mide qué guarda deja ver el rechazo, junto a un libro de envíos que descuenta de un tope declarado cada destinatario de una petición aceptada.
En MailerSend, la segunda ejecución del sondeo leyó la cuota al final, sin fiarse del contador de cada respuesta.
En Mailchimp Transactional no hay mitigación: por encima del umbral de la cola, el sondeo no mide.
Queda abierto qué pone al día el contador de MailerSend y cómo medir en Mailchimp Transactional por encima del umbral de la cola.

PRIME4API se ejecutó como servicio en local, no en el despliegue del ecosistema \citep{gavira2026tfm}.
Las 10 preguntas P de Mailchimp Transactional están calculadas a mano y sin contraste, porque el motor falla en 14 de sus 17 llamadas por dos huecos que el esquema admite.
El sesgo está en el instrumento: un despliegue con otra imagen podría responder de otro modo.
Lo mitiga que el servicio local sea la misma imagen de PRIME4API que despliega SPHERE en local (Sección~\ref{sec:hallazgos:limites-rates}).
No consta qué imagen ejecuta el despliegue del ecosistema.
En Mailchimp Transactional lo mitiga que el fallo esté declarado como fallo del instrumento, con los dos huecos que lo causan.
Queda abierto repetir las ejecuciones en el despliegue del ecosistema.

Los valores modelados dependen también de qué fuente se toma cuando dos discrepan.
La regla de la página de precios solo resuelve las divergencias en que la página es uno de los lados: 37 de las 82 registradas.
No resuelve las siete en que la página se contradice a sí misma.
Las demás siguen registrándose sin elegir.
Si la página publica un redondeo, el valor modelado es el redondeo, como el sobrecoste de Essentials 50K en SendGrid, que pasa de 0,00133 a 0,0013 USD.

Las Datasheets y el recuento de recurrencia los hicieron agentes de inteligencia artificial, y el agente que construye no es el que verifica (Anexo~\ref{anx:uso-ia}).
El riesgo es un dato que el agente dé por leído sin que la fuente lo contenga.
Lo mitigan la cita literal de cada valor, que una comprobación automática contrasta con la página archivada, y un agente verificador distinto del constructor, que vuelve a abrir cada fuente.
Queda abierto reunir lo que refutaron las verificaciones de las ocho APIs.

\subsection{Validez externa}\label{sec:conclusiones:externa}

La validez externa pregunta hasta dónde valen los resultados fuera del catálogo y de la toma de datos.
La respuesta a RQ2 se mide sobre ocho APIs de un solo dominio y con un solo consumidor de referencia (Sección~\ref{sec:evolucion:respuesta}).
Para otras APIs de correo vale el «no» de la primera pregunta, que no depende del tamaño de la muestra porque basta un caso que v0.3 no represente.
Al pasar de cuatro APIs a ocho, además, ninguna propuesta perdió su prueba.
Lo que depende de la muestra es la promoción, y no depende igual en todas las propuestas.
Es inferencia: las que recurren en más APIs, P-03 en las ocho, P-06 en seis y P-01 en cinco, son las que menos cambiarían con otro catálogo.
P-07, P-08 y P-09 descansan en una API cada una, así que su prueba no depende de la muestra y su alcance en el dominio queda sin medir.
Los cuatro elementos no soportados que cumplen el umbral con dos APIs justas son los que otro catálogo movería primero (Sección~\ref{sec:conclusiones:conclusion}).

P-04 y P-05 dependen del consumidor, no del dominio.
Es inferencia y no medida: cualquier Datasheet que declare un valor vacío o la cadena \texttt{custom} fallaría ante la misma revisión del motor.
Las demás propuestas nacen de casos de correo y no se han probado en otros dominios, así que la extensibilidad que anticipa el Capítulo~\ref{cap:introduccion} queda sin probar.
HARVEY4API no se ha ejecutado, y no hay un segundo consumidor que confirme los defectos de contrato.
Lo mitiga que cada propuesta lleve el caso que la motiva, de modo que otro dominio puede comprobarla caso a caso.
Queda abierto repetir el inventario en otro dominio y con un segundo consumidor.

La respuesta a RQ3 tiene una parte que no depende de la muestra.
Que la cuota sea la restricción activa sale del cociente, grande por construcción cuando el proveedor cobra por volumen mensual y publica un \textit{rate} por segundo o por minuto (Sección~\ref{sec:hallazgos:conclusiones}).
Las demás afirmaciones dependen de lo observado, y el paso de cuatro APIs a ocho debilitó tres de ellas.
El sesgo está en esas afirmaciones observadas, que otra muestra puede no repetir.
La comparación de RQ4 tiene el mismo límite: su base son ocho APIs y el estado de la suscripción desde el que se midió.
Para las afirmaciones observadas y para la comparación de RQ4 no hay mitigación en este trabajo.
Queda abierto repetir las afirmaciones observadas sobre otras APIs de correo.

La condición de representatividad de la Sección~\ref{sec:propuesta:respuesta} deja fuera de la prueba las formas de cobrar que no están en las ocho APIs.
Las cuatro formas salen del primer criterio de la consulta, que fijó el autor, y no de un censo del sector.
La respuesta registrada nombra además en la reserva una variante de cobro por créditos prepagados que el catálogo no incluye.
Lo mitiga que el criterio sea explícito y comprobable contra las páginas de precios, y que la consulta que completó la lista esté registrada.
Queda abierto si repetir la consulta daría la misma lista.

El catálogo modela solo el producto transaccional de cada proveedor, y el de \textit{marketing} queda fuera (Sección~\ref{sec:propuesta:seleccion}).
Las conclusiones sobre límites y \textit{rates} no se trasladan a ese producto, porque se cobra y se limita con otras unidades: Resend lo factura por contactos y no por correos enviados.
Lo mitiga que el ciclo de construcción no dependa del producto que se modela, así que el de \textit{marketing} puede recorrerlo igual.
Queda abierto construir sus Datasheets y comparar las dos respuestas.

La oferta de los proveedores cambia sin aviso, y el catálogo la describe tal como era en la ventana de toma de datos (Apartado~\ref{sec:propuesta:ventana}).
La propia selección lo muestra con el elemento no soportado de la fuente que caduca sin avisar (F-09): SES vendía planes antes de que la consulta lo describiera como pago por uso puro.
Un precio o un límite correcto en la toma de datos puede haber cambiado cuando se lea la memoria.
Lo mitiga que cada valor lleve su URL y su fecha de acceso (Sección~\ref{sec:propuesta:trazabilidad}).
Queda abierto un procedimiento de vigencia que diga cuándo reverificar cada Datasheet.

\subsection{Validez de conclusión}\label{sec:conclusiones:conclusion}

La validez de conclusión pregunta si los recuentos que sostienen las conclusiones están bien hechos y bien acotados.
El recuento de recurrencia deja 31 celdas sin verificar y 1 en suspenso (Tablas~\ref{tab:recurrencia} y~\ref{tab:recurrencia:b}).
Postmark y Resend reúnen 15 de esas celdas.
Una celda sin verificar solo puede sumar, así que cada recuento es un mínimo respecto de esas celdas, y en ellas el sesgo va en una sola dirección.
Verificarlas no retira ninguna promoción, pero podría hacer recurrentes a F-31, F-35 y F-36, y pasar F-16 de la traza a la recurrencia.
Confirmar en Mailgun que dos cuotas de la misma unidad se aplican a la vez haría recurrente F-38.
El umbral de dos APIs se cumple justo en cuatro filas, F-12, F-14, F-34 y F-37, así que en ellas la recurrencia descansa en dos celdas.
F-37 depende además de una regla provisional sobre qué variantes suman, y con otra regla bajaría a una sola API.
Lo mitiga que cada celda esté marcada fila a fila en las tablas del inventario.
Queda abierto verificar las 31 celdas y fijar la regla de variantes.

El recuento es del 23 de septiembre de 2026, y su fichero declara que las cuentas se comprobaron dos veces, a mano y con un script que las recalcula desde las celdas.
Ese script no está en el repositorio, así que la segunda comprobación no se puede repetir.
La doble comprobación acota los errores de aritmética, no el juicio sobre cada celda.
El 25 de septiembre de 2026 cada registro de elementos no soportados se alineó con el recuento en la columna de su propia API, y cada celda que cambió de lectura quedó anotada junto a la anterior.
Las copias de columnas de otras APIs que repiten las tablas acumulativas conservan su valor anterior.
En Mailgun, por ejemplo, la celda de F-09 pasó de presente a ausente, y la de F-07, de sin verificar a presente.
Siete celdas del recuento siguen sin fila en el registro de elementos no soportados de su API.
Ninguna suma a una promoción: cinco son ausencias, una es parcial y la séptima es una variante que tampoco cuenta.
Queda abierto versionar el script y crear esas siete filas.

Las coberturas del Capítulo~\ref{cap:hallazgos} dependen de dos criterios de recuento.
«No publicado» cuenta como respuesta cuando la pregunta lo admite, y en las preguntas de comparación una API cuenta cuando aporta su lado.
El primero no se aplica igual en todas las preguntas: en Postmark cuenta en la pregunta por el \textit{rate} máximo, y no en las que preguntan si hay límite diario o si el \textit{rate} sube con el plan.
El sesgo está en la cobertura de Postmark, que se movería en una o dos preguntas si la regla se unificara.
No hay mitigación en este trabajo: la regla se dejó como estaba y quedó anotada.
Queda abierto fijar esa regla en el catálogo de preguntas.

Las coberturas cuentan igual una respuesta en ámbar que una en verde.
En Postmark van en ámbar 13 de sus 18 respuestas y en Mailgun 12 de 20, frente a 1 de 22 en Resend y 1 de 20 en SendGrid.
El sesgo está en la comparación entre APIs: la cobertura dice cuánto se responde y no con qué confianza.
Lo mitiga que el Capítulo~\ref{cap:hallazgos} dé el reparto del ámbar por API.
Queda abierto dar la cobertura de cada API separada por color.

La otra convención aritmética es la del mes, que se toma de 30 días.
Con meses de 28 o de 31 días, los cocientes entre \textit{rate} y cuota mensual cambian un 6,7~\% o un 3,3~\%.
Con 28 días, el menor cociente de las cuentas sanas pasa de 2,6 a 2,4 y sigue por encima de uno.
La única fila por debajo de uno, MailerSend Starter 1M bajo revisión, sigue por debajo con los dos meses, así que ninguna conclusión se mueve.
Lo mitiga ese recálculo.
En esta convención no queda nada abierto.

\section{Paquete de replicación}\label{sec:conclusiones:replicacion}

El repositorio de datos que acompaña al trabajo contiene lo que sostiene los Capítulos~\ref{cap:catalogo} a~\ref{cap:evolucion}, y está disponible en \url{https://github.com/danvelcam/tfm-pricing-apis-correo}.
La Tabla~\ref{tab:conclusiones:paquete} recoge qué contiene cada ruta y qué parte de la memoria sostiene.

\begin{table}[H]
\centering
\small
\setlength{\tabcolsep}{3pt}
\begin{tabular}{|p{4.6cm}|p{6.2cm}|p{3.2cm}|}
\hline
\textbf{Ruta} & \textbf{Contenido} & \textbf{Sostiene} \\
\hline
\texttt{03\_datos/apis/} & Una carpeta por API: Datasheet, registro de fuentes, registro de elementos no soportados (\texttt{fricciones.md}), acta de revisión, matriz de \textit{features} y evidencias de los sondeos. Las copias fechadas de las páginas del proveedor no se publican: \texttt{PAGINAS.md} da la huella sha256 de cada una & Capítulo~\ref{cap:hallazgos}, Tablas~\ref{tab:recurrencia} y~\ref{tab:recurrencia:b} \\
\hline
\texttt{03\_datos/SELECCION\_\allowbreak APIS\_\allowbreak 2026-09-01.md} & Consulta de selección: prompt literal, respuesta íntegra, fecha y modelo & Sección~\ref{sec:propuesta:seleccion} \\
\hline
\texttt{03\_datos/RECURRENCIA\_\allowbreak 2026-09-23.md} & Recuento de recurrencia celda a celda & Tablas~\ref{tab:recurrencia} y~\ref{tab:recurrencia:b} \\
\hline
\texttt{03\_datos/CRITERIOS.md} & Reglas de modelado y de promoción & Capítulos~\ref{cap:catalogo} y~\ref{cap:evolucion} \\
\hline
\texttt{03\_datos/DISCREPANCIAS\_\allowbreak v03.md} & Discrepancias entre la guía v0.3, PRIME4API y el repositorio del ecosistema & Sección~\ref{sec:evolucion:compatibilidad} \\
\hline
\texttt{03\_datos/preguntas/} & Catálogo de preguntas de decisión & Sección~\ref{sec:hallazgos:catalogo-pi} \\
\hline
\texttt{03\_datos/\_plantilla/} & Plantilla de la Datasheet y lista de comprobación de la recolección & Sección~\ref{sec:propuesta:ciclo} \\
\hline
\texttt{05\_salida/\_cap5/} & Hechos por API, análisis y ejecuciones del motor & Capítulo~\ref{cap:hallazgos} \\
\hline
\texttt{verificar\_prime4api.py} & Herramienta de verificación contra PRIME4API & Tabla~\ref{tab:verificacion} \\
\hline
\texttt{gen\_datasheet\_viewer.py} & Validador de formato y estación de revisión & Sección~\ref{sec:evolucion:contrato} \\
\hline
\texttt{verificar\_citas.py} & Comprobación de cada cita literal de la Datasheet contra el volcado archivado de su página & Sección~\ref{sec:propuesta:trazabilidad} \\
\hline
\texttt{ejecutar\_prime4api\_PI.py} & Ejecución de las preguntas P contra PRIME4API & Tabla~\ref{tab:cobertura-pi} \\
\hline
\texttt{gen\_fig\_capacidad.py} & Curvas de capacidad con la biblioteca de PRIME4API & Figura~\ref{fig:capacidad} \\
\hline
\texttt{probe\_*}, \texttt{medir\_l2\_*.py} & Sondeos por API y comprobación de la condición de traza & Apartados de la Sección~\ref{sec:hallazgos:apis}, Sección~\ref{sec:evolucion:metodo} \\
\hline
\texttt{lib\_*} & Bibliotecas que cargan los sondeos: lectura de credenciales, guarda de destinatarios y tope de gasto & Apartados de la Sección~\ref{sec:hallazgos:apis} \\
\hline
\texttt{tests/} & Pruebas de las herramientas & --- \\
\hline
\end{tabular}
\caption{Contenido del repositorio de datos que acompaña al trabajo. Las herramientas, desde \texttt{verificar\_prime4api.py} hasta \texttt{tests/}, están en \texttt{05\_salida/\_herramientas/}. Sostiene: la parte de la memoria que se apoya en cada ruta.}
\label{tab:conclusiones:paquete}
\end{table}

La herramienta de verificación se ejecuta sobre una Datasheet con una sola orden desde la raíz del repositorio.

\begin{lstlisting}
python3 05_salida/_herramientas/verificar_prime4api.py \
    03_datos/apis/postmark/datasheet.yaml
\end{lstlisting}

La primera ejecución clona PRIME4API desde su repositorio público \citep{prime4apirepo}.
Necesita red y las dependencias de Python que lista la cabecera de la herramienta.
La orden devuelve 0 si pasan todas las comprobaciones y 1 si alguna falla.
Con la Datasheet de Postmark, lo esperado es la fila de la Tabla~\ref{tab:verificacion}: 11 comprobaciones superadas y 10 fallidas por valor ausente.
Esa fila se obtuvo con la revisión \texttt{8e6011c} del motor, cuyo último cambio es del 18 de agosto de 2026.
La herramienta clona la última revisión cuando no tiene copia, así que reproducir la fila exige situar antes la copia en esa revisión (Sección~\ref{sec:conclusiones:interna}).

El repositorio no incluye las credenciales de las APIs, que viven fuera de él.
Repetir un sondeo exige cuentas propias, y el resultado depende del estado de esas cuentas (Sección~\ref{sec:conclusiones:constructo}).
Tampoco incluye PRIME4API, que la herramienta clona al ejecutarse, ni las memorias de otros trabajos académicos citados en la bibliografía.
Sí incluye las copias fechadas que se archivaron de las páginas de los proveedores, que son la evidencia de las citas que el verificador comprueba (Sección~\ref{sec:conclusiones:constructo}).
De Postmark no se archivó ninguna.

\section{Trabajo futuro}\label{sec:conclusiones:futuro}

Los Capítulos~\ref{cap:hallazgos} y~\ref{cap:evolucion} dejan cinco líneas de investigación abiertas.
La primera es distinguir el límite publicado del aplicado (F-04).
SES declara una petición por segundo en sus acciones de gestión y respondió con \texttt{200} a 12 peticiones seguidas a 5,7 por segundo.
Mailchimp Transactional declara que el JSON de cada llamada debe ocupar menos de 10~MB, y respondió con \texttt{200} a una llamada de unos 40~MB.
El registro de elementos no soportados esboza la salida, que es distinguir en el esquema los dos límites con la fuente de cada uno, y queda por elegir qué cifra va al fichero (Sección~\ref{sec:evolucion:fuera}).
La segunda es fijar las 12 formas esbozadas de la Tabla~\ref{tab:sinforma} y contrastarlas con la guía.
La forma de la credencial que solo se gestiona desde la web (F-22) queda en la frontera del objeto del esquema, porque describe qué operaciones ofrece cada canal.
La tercera es el vínculo entre iPricing y la Datasheet, que daría instrumento a la pregunta funcional mixta.
Coincide con la línea «Pricing de servicios multicanal» que deja abierta Gavira Sánchez \citep{gavira2026tfm}, y este trabajo le aporta el caso del correo.
En ese caso, la multicanalidad se lee como el conjunto de interfaces de cada servicio, con la definición del NIST (Sección~\ref{sec:estado:acm}).
La cuarta es la medición formal, el cuarto peldaño de la escalera de validación \citep{datasheetusagetesting}, que mediría en una campaña de carga el 429 de Postmark y Mailgun y el caudal de SendGrid.
La quinta es ejecutar las preguntas H con HARVEY4API y repetir las preguntas P en el despliegue del ecosistema.

Dos líneas que Gavira Sánchez deja abiertas reciben de este trabajo un caso cada una \citep{gavira2026tfm}.
Completar la serialización ACM2YAML, que pide recoger los límites de petición individual, recibe el caso de los resultados por petición, documentados en siete APIs (F-05).
Validarla formalmente recibe la herramienta de verificación por uso.
Lo que este trabajo aporta a SPHERE (Capítulo~\ref{cap:estadoDelArte}) abre otras dos.
La primera es construir en la plataforma la sección de términos de consumo, con su pestaña de versiones y la procedencia de cada dato.
La segunda es asociar a cada interfaz sus límites de consumo, porque las \textit{skills} de Ruiz López ya toman la interfaz como unidad de análisis (Sección~\ref{sec:propuesta:criterios}).
Postmark recomienda diez conexiones por IP en el relé SMTP, pero no aplica ese tope.
Su API REST anuncia 5.000 peticiones por segundo en las cabeceras de respuesta.
El YAML del analizador nombra cada interfaz, pero no tiene un atributo para su límite.

Quedan además cuatro tareas del propio trabajo.

\begin{itemize}
\item Fijar la revisión del motor en la herramienta de verificación y registrarla en cada ejecución.
\item Verificar las 31 celdas del recuento y fijar la regla provisional de variantes.
\item Contrastar las propuestas con el código de la guía y con lo reportado en los repositorios del ecosistema después del 10 de agosto de 2026.
\item Definir un procedimiento de vigencia que diga cuándo reverificar cada Datasheet ante la fuente que caduca sin avisar (F-09).
\end{itemize}

Las demás amenazas de la Sección~\ref{sec:conclusiones:amenazas} dejan su parte abierta en su propio párrafo.

\section{Conclusiones técnicas}\label{sec:conclusiones:cierre}

\textbf{En seis de las ocho APIs, la cuota es la restricción activa.} El \textit{rate} sigue fijando el tiempo mínimo para enviar el volumen que la cuota deja (Tablas~\ref{tab:limites-rates} y~\ref{tab:limites-rates-b}), y en la cuenta de prueba de Mailchimp Transactional es lo único que lo fija.

\textbf{El estado de la suscripción y la reputación explican más de la capacidad observada que el plan.} En seis APIs hay al menos un estado publicado que condiciona la capacidad, y tres de ellas publican una cifra distinta por estado.
Un consumidor que compare planes sin mirar el estado compara cifras que su cuenta puede no tener.
P-06 da a ese eje el disparador y la declaración única que el mecanismo de \textit{aliases} de la guía no tiene.

\textbf{Con las ocho APIs, ACM2YAML v0.3 no basta, y las carencias de P-01 a P-03 están en dimensiones que ya modela.} El tope de carga expresado en otra métrica, que motiva P-03, aparece en las ocho APIs, y la dimensión que lo contiene, el \texttt{workload}, ya existe en v0.3.
La segunda observación del Capítulo~\ref{cap:estadoDelArte} partía de las dimensiones que v0.3 no serializa, como la concurrencia o los resultados por petición.
Los resultados por petición aparecen en siete APIs (F-05), y no llevan propuesta porque el esquema ya los reconoce como línea futura.
En las tres propuestas, el esquema declara con valores escalares dimensiones que el dominio vende por listas, consume por productos o acota en otra métrica.

\textbf{La conformidad con la guía no basta para que el motor consuma la Datasheet.} Todas las Datasheets del catálogo superan el validador de formato de este trabajo, y ninguna de las ocho es plenamente consumible.
En P-04 y P-05 el esquema admite el valor que describe al proveedor, y el consumidor de referencia falla al leerlo.
Por eso el método añade la verificación por uso a la de formato.

\textbf{Con las ocho Datasheets, el instrumental del canal API alcanza por clase 25 de las 29 preguntas no funcionales.} Cada Datasheet devuelve su propia cobertura, y esa diferencia entre APIs es el primer resultado de la comparación, antes de cualquier cifra de precio.
Con cuatro Datasheets había respuesta común en 17 preguntas.
Las cinco que se pierden con ocho caen por dos huecos de la página de precios de Brevo: ningún plan de pago documenta qué ocurre al superar la cuota, y los escalones intermedios de volumen no tienen precio.

\textbf{Con ocho APIs, el dominio muestra formas que las cuatro primeras no tenían.} Con cuatro APIs, el catálogo describía la parte regular del dominio.
Mailchimp Transactional no tiene planes ni tope de destinatarios en la API, y SES no tiene cuota incluida.
En tres APIs no hay sobrecoste.
La primera observación del Capítulo~\ref{cap:estadoDelArte} pedía confrontar el modelo con un conjunto representativo de las APIs de un sector, elegido con un criterio declarado, incluidas las que no encajan con naturalidad.
La respuesta es parcial en dos sentidos.
Las ocho APIs cubren cuatro formas de cobrar, no el sector entero (Sección~\ref{sec:conclusiones:externa}).
Y la Datasheet las describe con dos formas que la guía no define, mientras el motor no consume sin error ninguna de las ocho.
